% MI20091C
% Thoraxschmerz, beschwerdefrei
% 14.0
% 2013-11-30

:ref:`m-akutes-koronarsyndrom-beschwerdefrei`
%
\begin{itemize}
    -   Lagerung: Oberkörper hoch
  -   Bewegungensverbot, Schonung
  -   Beengende Kleidung öffnen
    -   O₂-Gabe gemäß \FullPrettyRef{m:sauerstoffberieselung}
  -   Bestmögliches Monitoring
  -   \II{Nitro-Spray}: Manche Patienten haben einen
  Nitroglyzerin-Spray (z. B. Nitrolingual{\texttrademark})
  zur Einnahme bei Symptomen eines Akuten Koronarsyndroms
  verschrieben bekommen. Der Spray soll weiter *vom
  Patienten selbstständig* in der *verschriebenen
  Dosierung* genommen werden, solange der systolische
  Blutdurck RR\textsubscript{syst} {\textgreater} 100\MmHg*
  ist. Der Helfer muss sich versichern, dass das Medikament
  für den jeweiligen Patienten verschrieben wurde und dass
  die vorgegebene Dosierung eingehalten wird. Eine
  regelmäßige Blutdruckkontrolle ist notwendig.
  -   Transportentscheidung: Abt. f. Innere Medizin
\end{itemize}